% !TeX root = ../../main.tex
\subsection{比例}

\begin{definitionbox}{\textbf{比例}}
    \label{def:比例}
    $y$ が $x$ の関数で，$x$ と $y$ の関係が次のような式で表されるとき、\textbf{$y$ は $x$ に比例する}という.

    \[\fitblankbf[]{y = ax} \qquad (a \neq 0)\]

    \vspace{0.5em}
    このとき、$a$ を\fitblankbf[]{傾き} あるいは \fitblankbf[]{変化の割合} という.${}^{*}$
\end{definitionbox}

比例は、一方の数量が2倍、3倍 … となると、それに伴ってもう一方の数量も2倍、3倍 … となるような関係をいう。

\begin{theorembox}{\textbf{比例のグラフ}}
    比例：$y = ax$ のグラフは\fitblankbf[]{原点を通る直線}である.
    \begin{enumerate}
        \item $a > 0$ のとき，グラフは\fitblankbf[]{右上がり}である.
        \item $a < 0$ のとき，グラフは\fitblankbf[]{右下がり}である.
    \end{enumerate}
\end{theorembox}

\begin{example}[【例題 1】]
    $y$ が $x$ に比例し，$x = 3$ のとき $y = 6$ である. $y$ を $x$ の式で表しなさい.
\end{example}

\begin{proof}\mbox{}\\
    $y = ax$ とおく.

    \begin{hiddenanswer}
        $x = 3$, $y = 6$ を代入すると，
        \begin{align*}
            6 &= 3a \\
            a &= 2
        \end{align*}
        \finalanswer{答え：$y = 2x$}
    \end{hiddenanswer}
\end{proof}

\vfill
\noindent\rule{\textwidth}{0.4pt}

\vspace{0.2em}
\noindent
$^{*}$ Definition \ref{def:比例} より、$x \neq 0$ のとき、$\bm{a = \frac{y}{x}}$ となる.

\newpage
